\documentclass[11pt]{amsart}
\usepackage[T1]{fontenc}
\usepackage[margin=1in]{geometry}
\usepackage{amsmath,amssymb}
\numberwithin{equation}{section}
\title{Notes on Discrete Energy Estimates}
\author{FlashTeX Fixtures Team}
\date{}
\begin{document}
\maketitle
\section{The basic inequality}
The central estimate of these notes is the discrete energy bound~\eqref{eq:energy} proved in Section~2, which controls the growth of solutions to the scheme~\eqref{eq:scheme} below. We consider sequences satisfying
\begin{equation}
u^{n+1} = u^{n} + \Delta t\, F(u^{n}),
\label{eq:scheme}
\end{equation}
where $F$ is Lipschitz with constant $L$. Summing the one-step bound gives the energy estimate
\begin{equation}
E^{n} \le e^{2Ln} E^{0} + C \Delta t,
\label{eq:energy}
\end{equation}
with $E^{n} = \|u^{n}\|^{2}$ and $C$ depending only on $F(0)$.
\section{Consequences}
Equation~\eqref{eq:scheme} is consistent of order one, and the stability bound~\eqref{eq:energy} from Section~1 converts consistency into convergence by the standard Lax argument. A useful variant sharpens the constant in~\eqref{eq:energy}:
\begin{equation}
E^{n} \le e^{Ln} E^{0} + C' (\Delta t)^{2},
\label{eq:sharp}
\end{equation}
and combining~\eqref{eq:sharp} with~\eqref{eq:scheme} yields convergence of order one. Finally, the triangle inequality together with~\eqref{eq:energy} and~\eqref{eq:sharp} gives continuous dependence on the data, which is exactly what Section~1 promised.
\end{document}
